\subsection{Paradigmen der datenbasierten Baumsegmentierung und -repräsentation}

Bei der Modellierung urbaner Baumstrukturen aus LiDAR-Daten erfordern insbesondere die Schritte der Einzelbaumsegmentierung und der anschließenden Visualisierung eine gezielte methodische Abwägung. Die jeweiligen Umsetzungen der beiden Phasen unterscheiden sich grundlegend hinsichtlich ihrer Datenstruktur, Recheneffizienz und ihres Detailgrads. Als theoretische Grundlage für spätere Designentscheidungen dieser Arbeit werden im Folgenden zunächst die raster- und punktwolkenbasierten Ansätze zur Einzelbaumsegmentierung vergleichend gegenübergestellt, bevor gängige Verfahren zur 3D-Repräsentation der extrahierten Baumstrukturen analysiert werden.